% =====================================================================
\chapter{Aproximación universal y profundidad}
\label{tema:3}
% =====================================================================

¿Qué funciones puede representar un perceptrón multicapa? Este tema demuestra que una sola capa oculta basta, en principio, para aproximar cualquier función continua en un compacto, y explica por qué ese resultado no justifica por sí solo el uso de redes anchas y poco profundas. La demostración en una dimensión es constructiva y usa solo la continuidad uniforme; el caso general se enuncia según Cybenko~\cite{cybenko}, Hornik~\cite{hornik} y Leshno \emph{et al.}~\cite{leshno}. La discusión sobre la profundidad sigue la §6.4.1 de~\cite{gbc} y el artículo de Telgarsky~\cite{telgarsky}; la §5.1.1.4 de~\cite{d2l} y el apartado «Representational power» de~\cite{cs231n} ofrecen el mismo punto de vista.

\section{Redes ReLU y funciones lineales a trozos}

\begin{definicion}[Función lineal a trozos]
Una función continua $f\colon\R\to\R$ es \emph{lineal a trozos} con \emph{codos} $t_1<\dots<t_m$ si es afín en cada uno de los intervalos $(-\infty,t_1]$, $[t_1,t_2]$, \dots, $[t_m,\infty)$. Cada uno de esos $m+1$ intervalos es una \emph{pieza}.
\end{definicion}

\begin{proposicion}\label{prop:piezas}
Sea $f(t)=c+\sum_{j=1}^{H}a_j\ReLU(w_jt+b_j)$ una red de una capa oculta con $H$ unidades ReLU, entrada y salida escalares. Entonces $f$ es continua y lineal a trozos con a lo sumo $H$ codos, es decir, con a lo sumo $H+1$ piezas.
\end{proposicion}

\begin{proof}
Si $w_j=0$, la unidad $j$ es constante. Si $w_j\neq 0$, la función $t\mapsto\ReLU(w_jt+b_j)$ es continua y afín a cada lado de $\tau_j=-b_j/w_j$. La suma de funciones continuas y afines a trozos es continua, y es afín en todo intervalo que no contenga en su interior ninguno de los puntos $\tau_j$. Hay a lo sumo $H$ puntos distintos.
\end{proof}

\begin{lema}[Sombrero]\label{lem:sombrero}
La función $\Lambda(t)=\ReLU(t)-2\ReLU(t-1)+\ReLU(t-2)$ vale $0$ para $t\le 0$ y para $t\ge 2$, vale $t$ en $[0,1]$ y $2-t$ en $[1,2]$.
\end{lema}

\begin{proof}
Para $t\le 0$ los tres términos se anulan. En $[0,1]$ solo el primero es no nulo y vale $t$. En $[1,2]$ se tiene $t-2(t-1)=2-t$. Para $t\ge 2$, $t-2(t-1)+(t-2)=0$.
\end{proof}

Sumando sombreros desplazados y escalados se obtiene cualquier poligonal con nodos equiespaciados; el teorema siguiente usa una construcción equivalente y más económica.

\section{El teorema en una dimensión}

\begin{teorema}[Aproximación universal en $\R$ con ReLU]\label{teo:au1d}
Sea $f\colon[a,b]\to\R$ continua y sea $\varepsilon>0$. Existe una red $g(t)=c+\sum_{j=1}^{H}a_j\ReLU(t-\tau_j)$ con una capa oculta tal que $\abs{f(t)-g(t)}<\varepsilon$ para todo $t\in[a,b]$.
\end{teorema}

\begin{proof}
\emph{Paso 1: poligonal.} Por el teorema de Heine, $f$ es uniformemente continua en $[a,b]$: existe $\eta>0$ tal que $\abs{f(s)-f(t)}<\varepsilon$ si $\abs{s-t}\le\eta$. Sea $n$ un entero con $(b-a)/n\le\eta$, sean $t_k=a+k(b-a)/n$ para $k=0,\dots,n$ y sea $I$ la poligonal que interpola $f$ en los nodos $t_k$. Si $t\in[t_k,t_{k+1}]$, entonces $I(t)=(1-\lambda)f(t_k)+\lambda f(t_{k+1})$ para algún $\lambda\in[0,1]$, y
\[
\abs{f(t)-I(t)}\le(1-\lambda)\abs{f(t)-f(t_k)}+\lambda\abs{f(t)-f(t_{k+1})}<\varepsilon,
\]
porque $t$ dista a lo sumo $\eta$ de ambos nodos.

\emph{Paso 2: la poligonal es una red.} Sea $m_k=(f(t_{k+1})-f(t_k))/(t_{k+1}-t_k)$ la pendiente de $I$ en $[t_k,t_{k+1}]$, y defínase
\[
g(t)=f(a)+m_0\ReLU(t-t_0)+\sum_{k=1}^{n-1}(m_k-m_{k-1})\ReLU(t-t_k).
\]
En $t=a=t_0$ todos los términos ReLU se anulan y $g(a)=f(a)=I(a)$. En el intervalo $(t_k,t_{k+1})$ están activas exactamente las unidades $0,\dots,k$, y la pendiente de $g$ es $m_0+\sum_{j=1}^{k}(m_j-m_{j-1})=m_k$, la misma que la de $I$. Dos funciones continuas en $[a,b]$ que coinciden en $a$ y tienen la misma pendiente en cada tramo son iguales; luego $g=I$ en $[a,b]$ y $\abs{f-g}<\varepsilon$. La red tiene $H=n$ unidades ocultas, con pesos de entrada~1, sesgos $-t_k$ y pesos de salida $m_0$ y $m_k-m_{k-1}$.
\end{proof}

\begin{ejemplo}[Aproximación de $t^{2}$]\label{ej:cuadrado}
En $[0,1]$ con $n=2$ tramos, los nodos son $0$, $1/2$ y $1$, con valores $0$, $1/4$ y $1$, y las pendientes son $m_0=1/2$ y $m_1=3/2$. La red es
\[
g(t)=\tfrac12\ReLU(t)+\ReLU\bigl(t-\tfrac12\bigr).
\]
En el primer tramo el error es $t/2-t^{2}$, máximo en $t=1/4$, donde vale $1/16$; en el segundo es $1/4+\tfrac32(t-\tfrac12)-t^{2}$, máximo en $t=3/4$, donde vale también $1/16$. Con $n$ tramos de longitud $h=1/n$ el error máximo es $h^{2}/4$ (\cref{eje:3-cuadrado}).
\end{ejemplo}

\section{El teorema general}

\begin{teorema}[Aproximación universal; Cybenko, Hornik, Leshno \emph{et al.}]\label{teo:au}
Sea $\phi\colon\R\to\R$ continua. Las funciones de la forma
\[
g(x)=\sum_{j=1}^{H}a_j\,\phi\bigl(x\cdot w_j+b_j\bigr),\qquad H\in\mathbb{N},\ a_j,b_j\in\R,\ w_j\in\R^{n},
\]
son densas en el espacio de funciones continuas sobre cualquier compacto $K\subset\R^{n}$ con la norma del supremo si y solo si $\phi$ no es un polinomio.
\end{teorema}

Cybenko~\cite{cybenko} lo demostró en 1989 para activaciones sigmoidales usando el teorema de Hahn--Banach y el de representación de Riesz; Hornik~\cite{hornik} lo extendió a activaciones acotadas y no constantes, y Leshno, Lin, Pinkus y Schocken~\cite{leshno} obtuvieron la caracterización anterior. La demostración del caso general excede el nivel de este curso. La necesidad de la condición sí es elemental:

\begin{proposicion}\label{prop:polinomio}
Si $\phi$ es un polinomio de grado $k$, toda red de la forma del \cref{teo:au} es un polinomio de grado a lo sumo $k$, y las redes no son densas en $C([0,1])$.
\end{proposicion}

\begin{proof}
Cada $\phi(x\cdot w_j+b_j)$ es un polinomio de grado $\le k$ en las componentes de $x$, y una combinación lineal de tales polinomios también lo es. El conjunto $\mathcal{P}_k$ de polinomios de grado $\le k$ en $[0,1]$ es un subespacio de dimensión finita de $C([0,1])$ y, por tanto, cerrado. Si las redes fueran densas, $C([0,1])\subseteq\overline{\mathcal{P}_k}=\mathcal{P}_k$, lo que es falso: $t^{k+1}\notin\mathcal{P}_k$.
\end{proof}

\section{Qué no dice el teorema: el papel de la profundidad}

El \cref{teo:au} es un resultado de existencia. No acota el número $H$ de unidades necesarias, que para algunas funciones crece exponencialmente con la dimensión $n$ de la entrada~\cite[§6.4.1]{gbc}; no garantiza que el descenso de gradiente encuentre los pesos, y no dice nada sobre la generalización a datos no vistos. La profundidad aporta una ventaja que el teorema no recoge: componer capas permite reutilizar piezas, y algunas funciones que una red profunda representa con pocas unidades requieren exponencialmente muchas en una red poco profunda. El ejemplo siguiente lo muestra de forma exacta.

\begin{proposicion}[Diente de sierra]\label{prop:sierra}
Sea $g(t)=2\ReLU(t)-4\ReLU(t-\tfrac12)$, que en $[0,1]$ vale $2t$ en $[0,\tfrac12]$ y $2-2t$ en $[\tfrac12,1]$, y sea $g^{(k)}=g\circ\cdots\circ g$ ($k$ veces). Entonces:
\begin{enumerate}
  \item en $[0,1]$, $g^{(k)}$ es lineal a trozos con exactamente $2^{k}$ piezas de longitud $2^{-k}$, y en cada una recorre todo $[0,1]$;
  \item $g^{(k)}$ se calcula con una red de $k$ capas ocultas de 2 unidades ReLU cada una;
  \item toda red de una capa oculta que coincida con $g^{(k)}$ en $[0,1]$ necesita al menos $2^{k}-1$ unidades.
\end{enumerate}
\end{proposicion}

\begin{proof}
(i) Por inducción sobre $k$. Para $k=1$ es la descripción de $g$. Supóngase cierto para $k$ y escríbase $g^{(k+1)}=g^{(k)}\circ g$. En $[0,\tfrac12]$, $g(t)=2t$ recorre $[0,1]$ de forma creciente y lineal, de modo que $g^{(k)}(2t)$ es una copia de $g^{(k)}$ comprimida en un factor~2: tiene $2^{k}$ piezas de longitud $2^{-k-1}$. En $[\tfrac12,1]$, $g(t)=2-2t$ recorre $[0,1]$ de forma decreciente y se obtiene otra copia, reflejada. En total hay $2^{k+1}$ piezas de longitud $2^{-(k+1)}$, y en cada una la función recorre $[0,1]$. (ii) Cada capa calcula $u\mapsto(\ReLU(u),\ReLU(u-\tfrac12))$; la combinación $2\ReLU(u)-4\ReLU(u-\tfrac12)$ que produce $g(u)$ se incorpora a los pesos de la capa siguiente, y la última se usa como salida. Como $g$ lleva $[0,1]$ en $[0,1]$, cada capa recibe valores en $[0,1]$. (iii) Por la \cref{prop:piezas}, una red con $H$ unidades tiene a lo sumo $H+1$ piezas; se necesita $H+1\ge 2^{k}$.
\end{proof}

Con $k=20$, la red profunda usa 40 unidades y la poco profunda más de un millón. Telgarsky~\cite{telgarsky} demostró que la separación persiste aunque solo se pida \emph{aproximar}: hay funciones calculables por redes de profundidad del orden de $k^{3}$ y anchura constante que una red de profundidad $O(k)$ solo aproxima con del orden de $2^{k}$ unidades. Los LLM son profundos por esta razón: LLaMA-7B apila 32 bloques~\cite[§2]{touvron}, cada uno con un MLP de anchura moderada.

\begin{lecturas}
\item Goodfellow, Bengio y Courville~\cite{gbc}, §6.4.1 «Universal Approximation Properties and Depth».
\item Zhang \emph{et al.}~\cite{d2l}, §5.1.1.4 «Universal Approximators».
\item Stanford CS231n~\cite{cs231n}, «Neural Networks Part~1», apartado «Representational power».
\item Cybenko~\cite{cybenko} y Leshno \emph{et al.}~\cite{leshno}, enunciados y demostraciones del caso general.
\item Telgarsky~\cite{telgarsky}, ventajas de la profundidad.
\end{lecturas}

\section*{Ejercicios del tema 3}
\addcontentsline{toc}{section}{Ejercicios del tema 3}

\begin{ejercicio}{1}{eje:3-existencia}
El \cref{teo:au} garantiza que existe una red de una capa oculta que aproxima cualquier función continua. Dé dos razones por las que no se entrenan LLM de una sola capa oculta enorme.
\end{ejercicio}

\begin{ejercicio}{2}{eje:3-codos}
Evalúe $g(t)=\ReLU(t)+2\ReLU(t-1)$ en $t=0$; $0{,}5$; $1$; $1{,}5$; $2$. ¿Qué pendientes tiene $g$ y dónde cambian?
\end{ejercicio}

\begin{ejercicio}{2}{eje:3-cuadrado}
(a) Construya, con el método del \cref{teo:au1d}, la red que interpola $f(t)=t^{2}$ en $[0,1]$ con $n=4$ tramos y calcule el error máximo. (b) Demuestre que con $n$ tramos de longitud $h$ el error máximo es $h^{2}/4$ y deduzca cuántas unidades bastan para un error menor que $10^{-4}$.
\end{ejercicio}

\begin{ejercicio}{2}{eje:3-sierra}
Con la función $g$ de la \cref{prop:sierra}, calcule $g^{(2)}(t)$ para $t=0,\tfrac18,\tfrac14,\dots,1$ y dibuje su gráfica. ¿Cuántas piezas tiene $g^{(3)}$ y cuántas unidades necesitaría una red de una capa oculta para representarla?
\end{ejercicio}

\begin{ejercicio}{3}{eje:3-abs-max}
Construya una red ReLU de una capa oculta que valga exactamente $\abs{t}$ en todo $\R$ y otra que valga $\max(t_1,t_2)$ en todo $\R^{2}$, sin término lineal directo de la entrada a la salida.
\end{ejercicio}

\begin{ejercicio}{3}{eje:3-cpwl}
Demuestre que toda función $\R^{n}\to\R$ calculada por un MLP con activaciones ReLU (en cualquier número de capas) es continua y lineal a trozos, en el sentido de que $\R^{n}$ se divide en un número finito de poliedros en cada uno de los cuales la función es afín.
\fuente{adaptado de Zhang \emph{et al.}~\cite{d2l}, §5.1, ejercicio~4.}
\end{ejercicio}
